\chapter{Signalling and Screening}\label{ch:b2:signalling-and-screening}

Earnings rise with schooling. If each year of school taught skills worth a little more pay, the rise would be gradual.
The chart shows average earnings in 1975 of American men by the number of years they spent at school. The rise is
not gradual: it jumps at twelve years, when a high-school diploma is earned, and at sixteen, when a college degree is.

\begin{openingchart}{Mean earnings in 1975 by years of schooling, in thousands of dollars: 15{,}992 men aged 16 to 55 of
  a 1970s United States population survey (the comparison sample of a job-training study). Dark bars: twelve and
  sixteen years. Earnings are capped at 25{,}244 dollars, the cap reached by 3{,}007 men; years 0 to 7 have fewer than
  300 men each.}{Dehejia and Wahba's NSW files (survey comparison sample); our computation}
\begin{tikzpicture}
\begin{axis}[omrecord, ybar=0pt, bar width=7pt, width=0.85\linewidth, height=0.42\linewidth, xmin=-0.7,
    xmax=18.7, ymin=0, ymax=18.5, xlabel={Years of schooling}, ylabel={Thousands of dollars},
    xtick={0,2,...,18}, minor xtick={1,3,...,17}, xmajorgrids=false]
  \addplot[fill=omThm!35, draw=omThm] table[x=education, y=other]
    {figdata/bachelor-2/16-signalling-and-screening/out/schooling.csv};
  \addplot[fill=omDef!80, draw=omDef] table[x=education, y=diploma]
    {figdata/bachelor-2/16-signalling-and-screening/out/schooling.csv};
\end{axis}
\end{tikzpicture}
\end{openingchart}

% figdata: summary (jump12 4.81, jump16 3.51, max_other 1.17, dollars11 9903, dollars12 14715, dollars15 12749, dollars16 16256)
Mean earnings rise from 9{,}903 dollars at eleven years to 14{,}715 at twelve, a jump of 4{,}812, and from 12{,}749 at
fifteen years to 16{,}256 at sixteen, a jump of 3{,}507. Between nine and eighteen years no other year adds more than
1{,}200 dollars, and several subtract. (Means at nine to eleven years are pulled down by young men still at school: the
sample starts at age 16.) A diploma seems to be worth more than the year of school that earns it.

One explanation is that the diploma is a \omterm{def:b2:signalling-and-screening:signal}{signal}. Employers cannot see a worker's productivity; they can see her
diploma; and finishing school is easier for those who will be productive. Then a diploma raises pay even if it adds no
skill. This chapter builds that model, then its mirror image, in which the uninformed side moves first and offers a
menu that makes the informed sort themselves. It ends with the evidence on how much of a diploma's value is \omterm{def:b2:signalling-and-screening:signal}{signal}.

\begin{recall}
From \cref{ch:b2:adverse-selection}: hidden information, \omterm{def:b2:adverse-selection:adverse-selection}{adverse selection} and the market for insurance. From
\cref{ch:b2:bayesian-games}: types, beliefs and \omterm{def:b2:bayesian-games:pbe}{perfect Bayesian equilibrium}. From the Year 1 volume: price
discrimination by a monopoly.
\end{recall}

\begin{notation}[Symbols of this chapter]\label{not:b2:signalling-and-screening:symbols}
In the signalling model, a worker's productivity is $\theta_L$ or $\theta_H$, her years of schooling $s$, and her cost of a year of
schooling $c_L$ or $c_H$; $s^*$ is the schooling level that the high type chooses. In \omterm{def:b2:signalling-and-screening:screening}{screening}, $q$ is a quantity bought,
$\mathit{F}$ a fee, $\theta_L$ and $\theta_H$ the consumer types, $n_L$ and $n_H$ their shares and $c$ the cost of a unit.
\end{notation}

\section{The job-market signalling model}\label{sec:b2:signalling-and-screening:signalling}

\begin{definition}[Signal]\label{def:b2:signalling-and-screening:signal}
A \emph{signal}\index{signal} is a costly action chosen by an informed party, which the uninformed party observes and
which is cheaper for some types than for others, so that it can reveal the type.
\end{definition}

Workers are of two types, equally many: low productivity $\theta_L = 10$ and high productivity $\theta_H = 20$, in thousands of
dalens a year. Before working, each chooses years of schooling $s$, which in this model add nothing to productivity. A
year costs the low type $c_L = 4$ and the high type $c_H = 2$, in effort, fees and earnings forgone: school is easier for
those who are more able. Employers see $s$ but not the type; competition among them sets the wage equal to expected
productivity given $s$. A worker's payoff is her wage minus $c\,s$.

The game is a signalling game: the worker moves first, the employers update their beliefs on seeing $s$, then pay. Its
equilibria are perfect Bayesian (\cref{def:b2:bayesian-games:pbe}).

\begin{definition}[Separating equilibrium]\label{def:b2:signalling-and-screening:separating}
A \emph{separating equilibrium}\index{separating equilibrium} is an equilibrium of a signalling game in which different
types choose different actions, so that the uninformed party learns each type from its action.
\end{definition}

\begin{definition}[Pooling equilibrium]\label{def:b2:signalling-and-screening:pooling}
A \emph{pooling equilibrium}\index{pooling equilibrium} is an equilibrium of a signalling game in which all types choose
the same action, so that the uninformed party learns nothing beyond its prior.
\end{definition}

Try this separating outcome: low types choose $s = 0$, high types $s = s^*$; employers believe that anyone with $s \ge s^*$ is high and
anyone else low, and pay 20 and 10. The low type does not imitate if $20 - 4s^* \le 10$; the high type does not prefer $s = 0$ if
$20 - 2s^* \ge 10$.

\begin{proposition}[The range of separating equilibria]\label{prop:b2:signalling-and-screening:range}
In the two-type signalling model, the outcome in which low types choose no schooling and high types choose $s^*$ is a
\omterm{def:b2:signalling-and-screening:separating}{separating equilibrium} if and only if
\[
  \frac{\theta_H - \theta_L}{c_L} \le s^* \le \frac{\theta_H - \theta_L}{c_H}.
\]
With the numbers above, $2.5 \le s^* \le 5$.
\end{proposition}

\begin{proof}
With beliefs as described, wages are $\theta_H$ for $s \ge s^*$ and $\theta_L$ otherwise. A worker choosing $s$ below $s^*$ gets $\theta_L$, so if she
does not choose $s^*$ she chooses $s = 0$; above $s^*$ she pays more for the same wage. So each type compares only 0 and $s^*$. The low
type keeps $s = 0$ if and only if $\theta_H - c_L s^* \le \theta_L$, which is the left inequality; the high type keeps $s^*$ if and only if
$\theta_H - c_H s^* \ge \theta_L$, the right one. Beliefs on the equilibrium path are correct by Bayes' rule, and off it any beliefs are
allowed, so these are all the conditions. The range is not empty because $c_H < c_L$.
\end{proof}

\begin{omfigure}
\begin{tikzpicture}[econ]
\begin{axis}[omrecord, grid=none, xmin=0, xmax=6.2, ymin=0, ymax=27, width=0.62\linewidth,
    height=0.45\linewidth, xlabel={Years of schooling $s$}, ylabel={Thousands of dalens a year},
    xtick={0,2.5,5}, ytick={10,15,20}, legend style={at={(1.03,0.5)}, anchor=west}]
  \addplot[black, very thick, const plot] coordinates {(0,10) (2.5,20) (6.2,20)};
  \addplot[omDef, thick, domain=0:4.25, samples=2] {10 + 4*x};
  \addplot[omThm, thick, dashed, domain=0:6.2, samples=2] {15 + 2*x};
  \node[eqpoint] at (axis cs:0,10) {};
  \node[eqpoint] at (axis cs:2.5,20) {};
  \node[anchor=north west, font=\footnotesize] at (axis cs:0,10) {low};
  \node[anchor=north west, font=\footnotesize] at (axis cs:2.5,20) {high};
  \legend{Wage schedule, Low type's indifference, High type's indifference}
\end{axis}
\end{tikzpicture}

{\small The least-cost \omterm{def:b2:signalling-and-screening:separating}{separating equilibrium}. The wage is 10 below $s^* = 2.5$ and 20 from it. The low type's indifference
line through her choice $(0, 10)$ rises by 4 a year and just reaches $(2.5, 20)$: imitating gains her nothing. The high type's
line rises by 2 a year; through $(2.5, 20)$ it lies above $(0, 10)$, so she prefers to \omterm{def:b2:signalling-and-screening:signal}{signal}.}
\end{omfigure}

\subsection{Pooling, and which equilibrium}

Pooling outcomes exist too. If both types choose $s_p$ and employers believe any other choice comes from a low type, the
wage at $s_p$ is the average, 15. The low type stays if $15 - 4s_p \ge 10$, that is $s_p \le 1.25$; the high type stays if $15 - 2s_p \ge 10$,
which is then implied. So every $s_p$ from 0 to 1.25 is a \omterm{def:b2:signalling-and-screening:pooling}{pooling equilibrium}.

The model has many equilibria because beliefs off the equilibrium path are free. A refinement called the
\emph{intuitive criterion} narrows them: an employer who sees a choice that one type could never gain from, whatever the
wage, should believe it was made by the other type. In the \omterm{def:b2:signalling-and-screening:separating}{separating equilibrium} with $s^* = 4$, a worker choosing $s = 3$
would get at most $20 - 12 = 8$ if low, less than her 10; so the choice must come from a high type, who gets $20 - 6 = 14$, more
than her $20 - 8 = 12$. So $s^* = 4$ fails, and by the same argument only the least-cost \omterm{def:b2:signalling-and-screening:separating}{separating equilibrium}, $s^* = 2.5$,
survives (\cref{exo:b2:signalling-and-screening:5}). In this example the pooling equilibria fail too.

\subsection{Signalling as waste}

In this model schooling produces nothing; it only sorts. In the least-cost \omterm{def:b2:signalling-and-screening:separating}{separating equilibrium} the high type earns 20
and spends 5 on schooling, keeping 15; the low type keeps 10. If signalling were impossible, both would be paid 15 and spend
nothing: the high type would be no worse off and the low type better off. With more high types the gap is starker: if 80
per cent were high, pooling would pay everyone 18, and the high type would gain from a ban on signalling too
(\cref{exo:b2:signalling-and-screening:3}). The private return to schooling exceeds its social return. Whether this matters
depends on how much of the return to schooling is \omterm{def:b2:signalling-and-screening:signal}{signal}, an empirical question (\cref{sec:b2:signalling-and-screening:diplomas}).

\begin{model}{Job-market signalling}
\assumes Productivity known to the worker only; schooling cheaper for the more productive; competitive employers who pay
expected productivity given what they see.
\predicts Pay that jumps at the levels of schooling that sort workers, such as diplomas, more than with years of schooling
as such; a private return to schooling above its social return.
\tested Jumps in earnings at diploma years in survey data (Hungerford and Solon 1987, as summarised by Clark and Martorell
2014; the opening chart).
\breaks When workers who barely pass and barely fail an exam that grants a diploma are compared, the diploma adds little
to earnings (Clark and Martorell 2014).
\end{model}

\section{Screening by a menu}\label{sec:b2:signalling-and-screening:screening}

\begin{definition}[Screening]\label{def:b2:signalling-and-screening:screening}
\emph{Screening}\index{screening} is the offer, by the uninformed party, of a menu of contracts designed so that each
type of informed party chooses a different one, revealing its type.
\end{definition}

\begin{definition}[Incentive-compatibility constraint]\label{def:b2:signalling-and-screening:ic}
An \emph{incentive-compatibility constraint}\index{incentive-compatibility constraint} requires that each type prefer
the contract designed for it to the contract designed for any other type.
\end{definition}

\begin{definition}[Participation constraint]\label{def:b2:signalling-and-screening:pc}
A \emph{participation constraint}\index{participation constraint} requires that each type prefer the contract designed
for it to not contracting at all.
\end{definition}

Take the insurance market of \cref{ch:b2:adverse-selection} with competing insurers that offer contracts rather than one
price. Everyone has wealth 100 (thousands of kests) and may lose 50; high risks lose with probability 0.5, low risks with
probability 0.2. Both have log utility. A contract is a cover $x$ bought at a premium; a person ends with $W_1$ if no loss
occurs and $W_2$ if it does. With fair premiums, a high risk buying cover $x$ pays $0.5x$, a low risk $0.2x$.

The separating menu gives high risks full cover at their fair premium, 25, so that $W_1 = W_2 = 75$. Low risks get cover at
their fair premium, but only up to the point where high risks do not want it: with log utility and probability 0.5, a
high risk is indifferent when $W_1 W_2 = 75^2 = 5{,}625$. With cover $x$ at the low-risk price, $(100 - 0.2x)(50 + 0.8x) = 5{,}625$,
or $0.16x^2 - 70x + 625 = 0$, whose smaller root is $x = 9.12$. Low risks get only 9.12 of cover against a loss of 50: the
price of separating them from high risks is that they are under-insured.

\begin{omfigure}
\begin{tikzpicture}[econ]
\begin{axis}[omrecord, grid=none, xmin=55, xmax=106, ymin=45, ymax=96, width=0.55\linewidth,
    height=0.55\linewidth, xlabel={$W_1$: wealth if no loss}, ylabel={$W_2$: wealth if loss},
    xtick={75,90,100}, ytick={50,57.29,75,90}, yticklabels={50,57.3,75,90},
    legend style={at={(1.05,0.5)}, anchor=west}]
  \addplot[black!50, thin, domain=45:96, samples=2] {x};
  \addplot[omDef, thick, domain=75:100, samples=2] {150 - x};
  \addplot[omThm, thick, domain=90:100, samples=2] {450 - 4*x};
  \addplot[omDef, thick, dashed, domain=60:100, samples=60] {5625/x};
  \addplot[omThm, thick, dashed, domain=93:103, samples=40] {exp((4.479 - 0.8*ln(x))/0.2)};
  \node[eqpoint] at (axis cs:100,50) {};
  \node[eqpoint] at (axis cs:75,75) {};
  \node[eqpoint] at (axis cs:98.18,57.29) {};
  \node[anchor=north east, font=\footnotesize] at (axis cs:99.5,49.5) {endowment};
  \node[anchor=south east, font=\footnotesize] at (axis cs:75,75) {high};
  \node[anchor=east, font=\footnotesize] at (axis cs:97.5,57.29) {low};
  \legend{Full cover, High risks' fair line, Low risks' fair line, High risks' indifference,
    Low risks' indifference}
\end{axis}
\end{tikzpicture}

{\small The separating menu. High risks get full cover on their fair line. Low risks get the point of their fair line on
the high risks' indifference curve through 75: any more cover and high risks would take it. The low risks' indifference
curve through their contract is steeper than the high risks' everywhere.}
\end{omfigure}

\begin{proposition}[No pooling equilibrium]\label{prop:b2:signalling-and-screening:no-pooling}
In a competitive insurance market with two risk types and insurers that offer contracts, no single contract bought by
both types at a break-even premium is an equilibrium.
\end{proposition}

\begin{proof}
A pooling contract that breaks even charges the average probability of loss, above the low risks' probability. At that
contract, a low risk will give up more wealth in the loss state for a unit of wealth in the no-loss state than a high risk
will, since a loss is less likely for her: her indifference curve is steeper. So there are contracts with slightly less
cover and a slightly lower premium that low risks prefer to the pooling contract and high risks do not. An insurer that
offers one attracts only low risks; since its premium is still close to the pooled one, above the low risks' expected
loss, it makes a profit. So the pooling contract cannot survive competition.
\end{proof}

The separating menu need not be an equilibrium either. If low risks are a large enough share, a single pooled contract
priced at the average probability, for instance full cover at the pooled premium, is better for them than their small
cover; high risks take it too, and an insurer that offers it breaks even. The separating menu is undercut, and no
equilibrium exists. The best such contract has partial cover, and it breaks the menu at a lower share than full cover does (\cref{exo:b2:signalling-and-screening:12}).

\section{Price discrimination as screening}\label{sec:b2:signalling-and-screening:pricing}

A seller with market power screens too. A phone company cannot tell light users from heavy users, but it can offer plans:
a small allowance for a small fee and a large one for a large fee. A user of type $\theta$ values $q$ units at $\theta q - q^2/2$; a unit
costs the company $c$. The light type, a share $n_L$, has $\theta_L$; the heavy type, a share $n_H$, has $\theta_H > \theta_L$.

\begin{definition}[Information rent]\label{def:b2:signalling-and-screening:information-rent}
An \emph{information rent}\index{information rent} is the surplus that a type keeps in a \omterm{def:b2:signalling-and-screening:screening}{screening} menu because of its
private information: what the party offering the menu must leave it so that it does not take another type's contract.
\end{definition}

With full information, the company would sell each type the quantity at which its marginal value equals cost, $\theta - c$, for
a fee equal to its whole value. Without it, a heavy user offered those two plans takes the light plan, which leaves him
surplus. The company must design the plans around four constraints: each type participates, and each prefers its own plan.

\begin{proposition}[The two-type screening menu]\label{prop:b2:signalling-and-screening:menu}
If the company offers two plans $(q_L, F_L)$ and $(q_H, F_H)$ to maximise profit, then (i) the heavy type gets the quantity of full
information, $q_H = \theta_H - c$ (no distortion at the top); (ii) the light type gets less than it would with full information,
\[
  q_L = \theta_L - c - \frac{n_H}{n_L}(\theta_H - \theta_L),
\]
or nothing when this is negative; (iii) the light type keeps no surplus, and the heavy type keeps an \omterm{def:b2:signalling-and-screening:information-rent}{information rent} of
$(\theta_H - \theta_L)\,q_L$.
\end{proposition}

\begin{proof}
At the optimum, the light type's \omterm{def:b2:signalling-and-screening:pc}{participation constraint} binds: otherwise both fees could be raised by the same amount.
The heavy type's incentive constraint binds too: otherwise $F_H$ could be raised. So $F_L = \theta_L q_L - q_L^2/2$, and the heavy user's
surplus equals what he would get from the light plan, $(\theta_H q_L - q_L^2/2) - F_L = (\theta_H - \theta_L)q_L$; that is his rent, and
$F_H = \theta_H q_H - q_H^2/2 - (\theta_H - \theta_L)q_L$. Profit is then
\[
  n_L\bigl(\theta_L q_L - \tfrac12 q_L^2 - c q_L\bigr) + n_H\bigl(\theta_H q_H - \tfrac12 q_H^2 - c q_H - (\theta_H - \theta_L) q_L\bigr).
\]
The first-order condition for $q_H$ gives $\theta_H - q_H = c$; that for $q_L$ gives $n_L(\theta_L - q_L - c) = n_H(\theta_H - \theta_L)$, which is (ii). The
two constraints left aside hold: the heavy type's surplus $(\theta_H - \theta_L)q_L$ is not negative, and the light type does not want
the heavy plan, since $q_H \ge q_L$ and the heavy plan charges for its extra units what the heavy type values them at, which is more
than the light type does.
\end{proof}

Each unit sold to light users raises the rent the company must leave heavy users, so it sells light users fewer units:
it gives up some surplus from them to extract more from heavy users. The listing finds the best menu by search over
both quantities, to check the formula.

\omcode{code/bachelor-2/16-signalling-and-screening/listings/menu.py}{1}{31}{The best two-plan menu, found by searching
over the two quantities.}\label{lst:b2:signalling-and-screening:menu}

\omoutput{code/bachelor-2/16-signalling-and-screening/outputs/menu.txt}{Light users get 2 units, not 4; heavy users get
8, as with full information, and keep a rent of 8.}\label{out:b2:signalling-and-screening:menu}

% code: listings/menu.py
With $\theta_L = 6$, $\theta_H = 10$, $c = 2$ and two thirds light users, the formula gives $q_L = 6 - 2 - \tfrac12 \times 4 = 2$ and $q_H = 8$, as the search
found. The light plan's fee, 10, is the light user's whole value of 2 units; the heavy plan's fee, 40, leaves the heavy
user $48 - 40 = 8$, exactly what he would get from the light plan.

\begin{omfigure}
\begin{tikzpicture}[econ]
\begin{axis}[omrecord, grid=none, xmin=0, xmax=10.5, ymin=0, ymax=52, width=0.62\linewidth,
    height=0.45\linewidth, xlabel={$q$: units in the plan}, ylabel={Fee},
    xtick={0,2,4,6,8,10}, ytick={10,20,40}, legend style={at={(1.03,0.5)}, anchor=west}]
  \addplot[omDef, thick, domain=0:10, samples=50] {6*x - x^2/2};
  \addplot[omThm, thick, dashed, domain=0.8:10, samples=50] {10*x - x^2/2 - 8};
  \addplot[black!60, thin, domain=0:10, samples=2] {2*x};
  \node[eqpoint] at (axis cs:2,10) {};
  \node[eqpoint] at (axis cs:8,40) {};
  \node[anchor=south east, font=\footnotesize] at (axis cs:2,10.5) {light plan};
  \node[anchor=south east, font=\footnotesize] at (axis cs:8,40.5) {heavy plan};
  \legend{Light user's value, Heavy user's value less 8, Cost}
\end{axis}
\end{tikzpicture}

{\small The \omterm{def:b2:signalling-and-screening:screening}{screening} menu. The light plan sits on the light user's value curve: he keeps nothing. Both plans sit on the
heavy user's indifference curve with surplus 8: he is indifferent between them and is assumed to take his own. Selling the
light user more units would move the light plan up the heavy user's value curve and raise his rent.}
\end{omfigure}

\begin{model}{Screening by a menu}
\assumes An uninformed party that moves first and can offer several contracts; buyers who choose the contract that is
best for them; a single trait on which buyers differ.
\predicts Each type sorts into its own contract; the type with the strongest demand gets the efficient contract and a rent;
the others get distorted contracts.
\tested Only consistency is reported here: UK annuity buyers of different mortality sorted into different products
(Finkelstein and Poterba 2004; \cref{ch:b2:adverse-selection}), but those products were priced on age and sex, not
designed as a menu. A test of the model needs a seller's menu and data on who chose what, to check the distortion at
the bottom and the absence of distortion at the top.
\breaks When buyers differ in several hidden traits at once, such as risk and aversion to risk, which can offset each other
(Chiappori and Salani\'e 2012).
\end{model}

\section{Diplomas: signal or skill?}\label{sec:b2:signalling-and-screening:diplomas}

\begin{definition}[Sheepskin effect]\label{def:b2:signalling-and-screening:sheepskin}
A \emph{sheepskin effect}\index{sheepskin effect} is the part of the earnings return to education that comes from
holding a diploma, over and above the return to the years of schooling it takes.
\end{definition}

The opening chart shows jumps at diploma years, but it compares different people. Those who finish twelve years are not
those who stop at eleven; they may differ in ability, health or family, and employers may see some of it. A jump in
raw means is consistent with signalling and with selection. To isolate the diploma, one needs people who differ in
holding it but in little else.

\begin{debate}[Diplomas: signal or skill?]
% ledger: b2:cm2014-design, b2:cm2014-results, b2:cm2014-prior, b2:cm2014-ged
Earlier survey studies, among them Hungerford and Solon (1987), estimated that diplomas and the completion of twelfth grade
come with large wage increases; in their own data Clark and Martorell (2014) find a diploma earnings differential of 10 to 20
per cent among workers who completed grade 12 and did not go to college. To separate \omterm{def:b2:signalling-and-screening:signal}{signal} from selection, Clark and
Martorell compared the earnings of workers who barely passed and barely failed the exit exam that a Texas student must pass
to earn a diploma. The diploma's estimated effect was always under 3 per cent of mean earnings and never statistically
different from zero; combined with Florida, the estimates rule out effects above 5 to 6 per cent. One reading: employers learn
productivity quickly and pay little for the paper itself. Another: this design measures the \omterm{def:b2:signalling-and-screening:signal}{signal} only for students at the
margin of passing, and for the high-school diploma only; for a General Educational Development certificate, Tyler, Murnane and
Willett (2000) estimated returns of 10 to 20 per cent for some workers, though later work found much smaller ones (both as
summarised by Clark and Martorell 2014).
\end{debate}

\section{Exercises}\label{sec:b2:signalling-and-screening:exercises}

\begin{exercise}[$\star$]\label{exo:b2:signalling-and-screening:1}
Read the opening chart. How large are the jumps at twelve and sixteen years, and how large is the largest other rise
between nine and eighteen years? Why might mean earnings at ten and eleven years be below those at eight?
\end{exercise}

\begin{exercise}[$\star$]\label{exo:b2:signalling-and-screening:2}
In the signalling model let $\theta_L = 12$, $\theta_H = 18$, $c_L = 3$ and $c_H = 1.5$. Find the range of separating equilibria.
\end{exercise}

\begin{exercise}[$\star$]\label{exo:b2:signalling-and-screening:3}
Predict, then compute: in the model of the chapter, compare each type's payoff in the least-cost \omterm{def:b2:signalling-and-screening:separating}{separating equilibrium}
with pooling at $s = 0$, first with half the workers high, then with 80 per cent.
\end{exercise}

\begin{exercise}[$\star$]\label{exo:b2:signalling-and-screening:4}
For the menu found by the listing, compute each user's surplus from his own plan and from the other, and check the four
constraints.
\end{exercise}

\begin{exercise}[$\star\star$]\label{exo:b2:signalling-and-screening:5}
In the \omterm{def:b2:signalling-and-screening:separating}{separating equilibrium} with $s^* = 4$, show that a deviation to $s = 3$ could only come from a high type, and that it pays.
Explain why every $s^*$ above 2.5 fails in the same way.
\end{exercise}

\begin{exercise}[$\star\star$]\label{exo:b2:signalling-and-screening:6}
With 80 per cent high types, find the range of pooling equilibria in which employers believe any other choice comes from a
low type.
\end{exercise}

\begin{exercise}[$\star\star$]\label{exo:b2:signalling-and-screening:7}
Check the low-risk contract of \cref{sec:b2:signalling-and-screening:screening}: compute $W_1$ and $W_2$, the premium, and the low
risk's expected utility. Does she prefer it to the high risks' contract?
\end{exercise}

\begin{exercise}[$\star\star$]\label{exo:b2:signalling-and-screening:8}
Rerun the \omterm{def:b2:signalling-and-screening:screening}{screening} model of the listing with half the users light. What does the company do with light users, and what is
its profit per user?
\end{exercise}

\begin{exercise}[$\star\star$]\label{exo:b2:signalling-and-screening:9}
What does the human-capital view, in which each year of school adds skill, predict for the opening chart, and what does
signalling predict? Why does the chart alone not settle the question?
\end{exercise}

\begin{exercise}[$\star\star\star$]\label{exo:b2:signalling-and-screening:10}
Show that the \omterm{def:b2:signalling-and-screening:separating}{separating equilibrium} with $s^* = (\theta_H - \theta_L)/c_L$ is the only one that survives the intuitive criterion.
\end{exercise}

\begin{exercise}[$\star\star\star$]\label{exo:b2:signalling-and-screening:11}
In \cref{prop:b2:signalling-and-screening:menu}, find the share of heavy users above which the company sells nothing to light
users. Evaluate it for the listing's numbers.
\end{exercise}

\begin{exercise}[$\star\star\star$]\label{exo:b2:signalling-and-screening:12}
In the insurance market of \cref{sec:b2:signalling-and-screening:screening}, let 90 per cent of people be low risks. Find the
pooled fair premium for full cover and the low risks' utility from it. Is the separating menu an equilibrium? Find the
share of low risks above which it is not.
\end{exercise}

\section{Problem: Brenmoor phone plans}\label{sec:b2:signalling-and-screening:problem}

\begin{problem}[{Weekend problem --- a phone company in Brenmoor cannot tell light users from heavy users: what plans does it
offer, who keeps a rent, and what does the screening cost in surplus}]
\label{pb:b2:signalling-and-screening:1}
The company has 9{,}000 customers; 6{,}000 are light users and 3{,}000 heavy. A user of type $\theta$ values $q$ gigabytes a month at
$\theta q - q^2/2$ dalens, with $\theta = 30$ for light and $\theta = 40$ for heavy users. A gigabyte costs the company 10 dalens.

\medskip
\noindent\textbf{Part I --- Two consumer types.}
\begin{enumerate}
  \item Write each type's marginal value of a gigabyte.
  \item How many gigabytes would each type buy at a price of 10 a gigabyte?
  \item What is each type's surplus at that price?
  \item Which type values any given allowance more, and by how much per gigabyte?
  \item What does the company know, and what does it not?
\end{enumerate}

\medskip
\noindent\textbf{Part II --- Full-information plans.}
\begin{enumerate}[resume]
  \item If the company could tell the types apart, what allowance and fee would it give each?
  \item What is its profit per light and per heavy user, and in total?
  \item What surplus do users keep?
  \item If it offered the two plans openly, which would a heavy user take, and with what surplus?
  \item What would the company's profit then be?
  \item Would it do better selling only the heavy plan?
\end{enumerate}

\medskip
\noindent\textbf{Part III --- The screening menu.}
\begin{enumerate}[resume]
  \item Write the two participation and two incentive constraints.
  \item Which two bind at the optimum, and why?
  \item Find the light plan's allowance.
  \item Find the heavy plan's allowance.
  \item Find the light plan's fee.
  \item Find the heavy user's rent and the heavy plan's fee.
  \item Check that the light user does not want the heavy plan.
  \item Find the profit per user of each type and in total.
\end{enumerate}

\medskip
\noindent\textbf{Part IV --- Welfare and the information rent.}
\begin{enumerate}[resume]
  \item What is the total of the heavy users' rents?
  \item How much surplus is lost on each light user, and in total?
  \item Show that the company's lost profit, compared with full information, equals the rents plus the lost surplus.
  \item If half the customers were heavy users, what allowance would light users get?
  \item Above what share of heavy users would light users get nothing?
  \item State the result: the menu, the distortion at the bottom and the heavy users' rent, in dalens.
\end{enumerate}
\end{problem}
